% TOP_PROBLEM: {"id": 20000582, "title": "Finite-dimensional embeddings of doubling Hilbert subsets", "category_id": 8, "category": {"id": 8, "name": "analysis", "display_name": "Analysis", "description": "Limits, continuity, calculus, and function theory.", "slug": "analysis", "order_index": 8, "created_at": "2026-07-31T15:26:25.671Z"}, "status": "partially_solved"}
% FIRST_POSTED: 2026-10-06
% ATTEMPT_STATUS: solved
% OPENAI_CLAIM: {"schema_version":1,"source_repo":"https://github.com/openai/math","source_commit":"adc7f1241b42e322a6451854ab7e4b4c146bf78a","release_date":"2026-10-06","match":"full","resolution":"disproved","independently_reviewed":false,"second_pass":"agreed","families":[{"family":"098","title":"Compact counterexamples to bi-Lipschitz dimension reduction","subject":"Convex and metric geometry","manuscripts":[{"title":"A doubling Hilbert subset with no finite-dimensional bi-Lipschitz embedding","pdf_path":"preprints/A-doubling-Hilbert-subset-with-no-finite-dimensional-bi-Lipschitz-embedding-September-25-2026/main.pdf","readme_path":"preprints/A-doubling-Hilbert-subset-with-no-finite-dimensional-bi-Lipschitz-embedding-September-25-2026/README.md","date":"2026-09-25","theorem_ref":"thm:main (main.tex:47)","statement_sources":[{"path":"preprints/A-doubling-Hilbert-subset-with-no-finite-dimensional-bi-Lipschitz-embedding-September-25-2026/build/main.tex","line":47,"label":"thm:main"}]}],"lean":{"doc_path":"lean/docs/098.md","comparators":["lean/ComparatorChallenges/CompactBanach.json","lean/ComparatorChallenges/CompactBanach.lean","lean/ComparatorChallenges/DoublingHilbert.json","lean/ComparatorChallenges/DoublingHilbert.lean"],"declarations":["OAI.CompactBanach.main","OAI.DoublingHilbert.main"],"covers_main_theorem":true},"reasoning_trace":null}]}
\documentclass[11pt]{article}
\usepackage[margin=25mm]{geometry}
\usepackage{fontspec}
\setmainfont{Latin Modern Roman}
\usepackage{amsmath,amssymb,mathtools}
\usepackage{unicode-math}
\setmathfont{Latin Modern Math}
\usepackage{enumitem}
\usepackage{xurl,hyperref,bookmark}
\hypersetup{colorlinks=true,urlcolor=blue}
\setcounter{secnumdepth}{0}
\setlength{\parindent}{0pt}
\setlength{\parskip}{5pt}
\setlength{\emergencystretch}{3em}
\begin{document}
\section*{Finite-dimensional embeddings of doubling Hilbert subsets}

\subsection{OpenAI's claim}
OpenAI presents the following claims relevant to this question, with manuscript numbers following paper order in Sources. In manuscript 1 (family 098; 2026-09-25), the asserted result is that in each infinite-dimensional real Banach space one can choose a compact subset with a universal doubling constant for which every bi-Lipschitz map into a finite-dimensional normed target is impossible.

\subsection{Mathematical statement of the claimed result}
\paragraph{Manuscript 1.}
A single subset $S$ of real $\ell_2$, with the distance inherited from its Hilbert norm, is claimed to satisfy
\[
 \operatorname{dbl}(S)\le76800,\qquad
 \forall k\ge1\ \forall D<\infty:\ S\text{ has no bi-Lipschitz embedding into }\mathbb R^k\text{ of distortion }D.
\]
Here doubling means that every metric ball is covered by at most that many balls of half its radius. The obstruction concerns one fixed doubling space, not a sequence with increasing doubling constants.

\subsection{Outline of the argument}
For manuscript 1 (family 098; 2026-09-25), sheet-like pieces are arranged so that coverings remain uniformly doubling. Near-extremal derivatives of a hypothetical finite-dimensional embedding have small quadratic variation; crossings between separated sheets then force more packed images than the target can accommodate.

\subsection{Scope relative to this problem}
The displayed question has a negative answer under the manuscript's stated construction. Its counterexample meets the conditions in this problem and defeats the universal assertion, so this is a full disproof claim. The supporting result is manuscript 1 (family 098; 2026-09-25).

\subsection{Formal verification}
For manuscript 1 (family 098; 2026-09-25), OpenAI lists Lean statements for this assertion: in each infinite-dimensional real Banach space one can choose a compact subset with a universal doubling constant for which every bi-Lipschitz map into a finite-dimensional normed target is impossible.

The selected statement includes the clause used for this problem.

The linked family Comparator checks are \texttt{CompactBanach.lean}, \texttt{DoublingHilbert.lean}.

Their named declarations include \texttt{OAI.CompactBanach.main}, \texttt{OAI.DoublingHilbert.main}.

This site has not executed the Lean code. Described statement coverage is not an independent proof check.

\subsection{Status caveat}
The manuscripts were produced by an internal OpenAI model. OpenAI cautions that results without formalization can contain mistakes. This site has not independently checked the proofs; the label records the scope of OpenAI's claim.

\subsection{Sources}
\begin{itemize}
\item OpenAI, \emph{A doubling Hilbert subset with no finite-dimensional bi-Lipschitz embedding}, 2026-09-25; citation key \texttt{OAI:A-doubling-Hilbert-subset-with-no-finite-dimensional-bi-Lipschitz-embedding-September-25-2026}. \href{https://github.com/openai/math/blob/adc7f1241b42e322a6451854ab7e4b4c146bf78a/preprints/A-doubling-Hilbert-subset-with-no-finite-dimensional-bi-Lipschitz-embedding-September-25-2026/main.pdf}{PDF (GitHub)}; \href{https://github.com/openai/math/raw/adc7f1241b42e322a6451854ab7e4b4c146bf78a/preprints/A-doubling-Hilbert-subset-with-no-finite-dimensional-bi-Lipschitz-embedding-September-25-2026/main.pdf}{PDF (direct)}; \href{https://github.com/openai/math/blob/adc7f1241b42e322a6451854ab7e4b4c146bf78a/preprints/A-doubling-Hilbert-subset-with-no-finite-dimensional-bi-Lipschitz-embedding-September-25-2026/README.md}{Manuscript page and citation}; \href{https://github.com/openai/math/blob/main/preprints/A-doubling-Hilbert-subset-with-no-finite-dimensional-bi-Lipschitz-embedding-September-25-2026/main.pdf}{Latest version}; Statement source: \href{https://github.com/openai/math/blob/adc7f1241b42e322a6451854ab7e4b4c146bf78a/preprints/A-doubling-Hilbert-subset-with-no-finite-dimensional-bi-Lipschitz-embedding-September-25-2026/build/main.tex\#L47}{thm:main}.
\item Family 098: \href{https://github.com/openai/math/blob/adc7f1241b42e322a6451854ab7e4b4c146bf78a/lean/docs/098.md}{Lean scope}.
\item \href{https://github.com/openai/math/blob/adc7f1241b42e322a6451854ab7e4b4c146bf78a/lean/ComparatorChallenges/CompactBanach.json}{Comparator CompactBanach.json}.
\item \href{https://github.com/openai/math/blob/adc7f1241b42e322a6451854ab7e4b4c146bf78a/lean/ComparatorChallenges/CompactBanach.lean}{Comparator CompactBanach.lean}.
\item \href{https://github.com/openai/math/blob/adc7f1241b42e322a6451854ab7e4b4c146bf78a/lean/ComparatorChallenges/DoublingHilbert.json}{Comparator DoublingHilbert.json}.
\item \href{https://github.com/openai/math/blob/adc7f1241b42e322a6451854ab7e4b4c146bf78a/lean/ComparatorChallenges/DoublingHilbert.lean}{Comparator DoublingHilbert.lean}.
\item \href{https://github.com/openai/math/blob/adc7f1241b42e322a6451854ab7e4b4c146bf78a/CONTENTS.md}{OpenAI catalogue entry}.
\end{itemize}
\end{document}
